% !TEX root = ../main.tex
\section{엔도스랭크와 AWP를 나란히 놓고 보기}
\label{ch:appendix-er-awp}

\ref{ch:results}장은 이 논문의 모든 점수를 함께 보고해요. 그리고 각 페이지랭크(PageRank)를, 그 페이지랭크가 매끄럽게 다듬는(평활하는) 원시 차수(화살표를 그냥 센 수)와 견주어 읽어요. 이 부록은 한 층만 걷는 두 점수인 엔도스랭크(EndorseRank)와 AWP(적응형 가중 페이지랭크)를 바로 견준 결과를 모아요. 같은 기간에서의 지표 가족(비슷한 지표 묶음) 평균과 짝지은 차이, 대리 지표(대신 재는 값) 하나하나에 대한 계수, 그리고 기간 밖에서의 짝지은 차이예요. 표~\ref{tab:endorserank-awp-holdout}의 트레이더(거래하는 사람) 행과 끝에 있는 격자 표를 빼면, 이 비교는 어느 것도 미리 정한 것이 아니고, 어느 것도 이 논문의 주장(C)을 담고 있지 않아요. 그 트레이더 행과 격자 표의 쌍은 W1과 함께 등록했고, 여기서는 재분석으로 고친 점수에서 계산해요. 두 점수는 서로 다른 화살표(엣지) 위를 걷고, 서로 다른 구성개념(점수가 재려고 하는 생각)을 위해 만들었어요. 그래서 점검마다, 앞서는 점수는 거의 언제나 그 점검이 세는 화살표를 걷는 점수예요. 두 화살표 종류 어느 쪽으로도 만들지 않은 결과는 예외예요. 어느 점수도 그 결과를 위해 만들지 않았기 때문이에요. 그 결과는 \ref{sub:neutral-label-results}절에서 보고해요. 두 점수가 점수를 매기는 범위도 달라요. 매칭 코호트(모든 점수가 똑같이 점수를 매기는 계약 묶음)의 계약 가운데 $191$개에는 허락 화살표가 하나도 닿지 않고, $1{,}192$개에는 송금 화살표가 하나도 닿지 않아요. 기간 밖에서는 어느 점수도 자기가 다듬는 원시 차수를 이기지 못해요. 나중의 새 송금자(새로 송금을 보내온 주소)를 줄 세우는 데서 AWP는 송금 입차수(송금을 보내온 서로 다른 주소 수)보다 못해요. 그리고 나중의 새 허락 쌍을 줄 세우는 데서 엔도스랭크는 받은 허락 수(허락 입차수)보다 못해요. 고르게 다시 출발할 때도, 활동 재시작(활발한 주소에서 더 자주 다시 출발하기)을 쓸 때도 그래요(\ref{sec:holdout-results}절과 \ref{sub:endorserank-restarts}절; 라벨을 본 뒤에 계산했어요). 그러므로 엔도스랭크가 AWP보다 앞선다고 해서, 허락을 따라 퍼뜨리는 것(전파)이 허락한 주인 수에 무엇인가를 더한다는 것이 드러나지는 않아요.

\dissertationSubheading[sec:er-awp-same-window]{같은 기간}

표~\ref{tab:alignment-hybrid}은 두 점수의 지표 가족 평균을 구간(믿을 만한 범위)과 함께 보여 줘요. 표~\ref{tab:tie-sensitivity}의 마지막 행은 두 점수 사이의 일치를 보여 줘요. 두 모양은 서로 엇갈려요. AWP는 송금 가족을 가장 가깝게 따라가고, 허락 가족을 가장 덜 따라가요. 반면에 엔도스랭크는 허락 가족을 가장 가깝게 따라가요. 표~\ref{tab:tau-diff}은 계약을 다시 뽑은 같은 표본에서 짝지은 차이 다섯 개를 검사해요(\ref{sec:statistical-tests}절). 이 차이들에는 표의 행 순서대로 i부터 v까지 번호를 붙였어요.

% Generated by experiments/scripts/export_latex.py -- do not edit by hand.
% Source: experiments/data/2-processed-tables-and-evaluations/revised-usd-12m/same-window/eval_summary.json
\begin{table}[htbp]
\centering
\caption[같은 기간에서 엔도스랭크와 AWP 사이의 짝지은 대비.]{같은 기간에서 엔도스랭크와 AWP의 켄달의 $\tau$ 계수 사이 차이를 짝지은 부트스트랩으로 검사한 결과예요(매칭 코호트, 계약 $n=6{,}483$개). 지표 가족 항목에는 지표 가족 평균 $\tau$를 써요. $a$와 $b$에 같은 계약 다시 뽑기를 써요. 그래서 구간은 짝지은 차이에 대한 구간이에요. 구간이 0을 벗어나면, 5\% 수준에서 표본을 뽑을 때 생기는 흔들림만으로는 설명할 수 없는 차이라는 뜻이에요. 신뢰구간(CI)은 계약을 다시 뽑아 구한 95\% 백분위 부트스트랩 구간이에요(짝지은 다시 뽑기 400번, 시드 42).}
\label{tab:tau-diff}
\fitwidth{%
\begin{tabular}{lccccc}
\toprule
대비 ($a$ 빼기 $b$) & $\tau(a)$ & $\tau(b)$ & $\Delta\tau$ & $\Delta\tau$의 95\% 신뢰구간 & $P(\Delta\tau>0)$ \\
\midrule
EndorseRank 허락 빼기 EndorseRank 송금 & 0.560 & 0.196 & +0.363 & [0.341, 0.382] & 1.000 \\
AWP 송금 빼기 EndorseRank 송금 & 0.736 & 0.196 & +0.539 & [0.518, 0.560] & 1.000 \\
EndorseRank 허락 빼기 AWP 허락 & 0.560 & 0.226 & +0.333 & [0.311, 0.354] & 1.000 \\
AWP 시빌 안정성 빼기 EndorseRank 시빌 안정성 & 0.419 & 0.026 & +0.393 & [0.373, 0.412] & 1.000 \\
EndorseRank 받은 허락 수 빼기 EndorseRank 입차수 & 0.650 & 0.218 & +0.432 & [0.410, 0.452] & 1.000 \\
\bottomrule
\end{tabular}
}
\end{table}


다섯 구간은 모두 0보다 위에 있어요. AWP의 송금 일치도는 엔도스랭크의 송금 일치도보다 높아요(대비~ii, $+0.539$, $[0.518,0.560]$). 엔도스랭크의 허락 일치도는 AWP의 허락 일치도보다 높아요(대비~iii, $+0.333$, $[0.311,0.354]$). 그리고 AWP는 시빌 안정성 가족(시빌 공격, 곧 가짜 주소를 많이 만들어 속이는 공격을 생각해 고른 지표 묶음)에서 엔도스랭크보다 앞서요(대비~iv, $+0.393$, $[0.373,0.412]$). 나머지 두 대비는 엔도스랭크를 자기 자신과 견줘요. 엔도스랭크의 허락 일치도는 자기의 송금 일치도보다 높아요(대비~i, $+0.363$, $[0.341,0.382]$). 그리고 받은 허락 수와의 일치도는 들어오는 화살표 수(입차수)와의 일치도보다 높아요(대비~v, $+0.432$, $[0.410,0.452]$). 엔도스랭크는 계약 쌍의 $0.20\%$만 동점으로 둬요. 그래서 엔도스랭크 자신의 동점은 이 계수들 어느 것에도 천장(얻을 수 있는 가장 큰 값)을 씌우지 않아요. 그리고 대리 지표의 동점으로는 이만큼 큰 차이를 설명할 수 없어요(부록~\ref{sec:kendall-properties}). 허락 그래프 밖의 계약들에 0점을 주는 대신 외톨이 점(화살표가 하나도 없는 점)으로 남겨 두어도, 지표 가족 평균은 하나도 바뀌지 않아요(표~\ref{tab:tie-sensitivity}).

대비~ii와~iii의 방향이 서로 반대인 것은 구성개념에서 나와요. 각 점수는 자기 화살표로 만든 가족에서 앞서요. 그리고 가족마다 주된 대리 지표에는, 그에 맞는 페이지랭크가 다듬는 입차수가 들어 있어요. 그래서 이 일치는 일부가 만든 방식 때문에 저절로 생기고, 각 점수의 의도한 구성개념을 점검해요. 어느 차이도 어느 점수가 평판을 더 잘 재는지는 말해 주지 않아요.

\dissertationSubheading[sec:proxy-family-detail]{대리 지표 하나하나 보기}

지표 가족 평균 하나는 대리 지표마다 구한 계수 여러 개를 한데 줄인 값이에요. 표~\ref{tab:alignment-transfer}부터~\ref{tab:alignment-sybil-stability}까지는 엔도스랭크와 AWP의 계수를 대리 지표 하나씩 보고해요. 스피어먼의 로 $\rho$(두 줄 세우기가 얼마나 비슷한지 나타내는 수), 켄달의 타우 $\tau$(두 줄 세우기가 얼마나 같은 순서인지 나타내는 수), 그리고 $\tau$의 95\% 구간이에요. 송금 대리 지표는 AWP를 만든 사람들이 AWP를 평가할 때 쓴 바로 그 지표예요\parencite{do2023awp}.

\dissertationSubsubheading[sub:transfer-family]{송금 가족}

표~\ref{tab:alignment-transfer}은 두 점수와 입차수, 그리고 받은 양의 합 사이의 계수를 보여 줘요.

% Generated by experiments/scripts/export_latex.py -- do not edit by hand.
% Source: experiments/data/2-processed-tables-and-evaluations/revised-usd-12m/same-window/eval_summary.json
\begin{table}[htbp]
\centering
\caption[송금 대리 지표에 대한 엔도스랭크와 AWP의 분야 일치도.]{엔도스랭크와 AWP가 송금 대리 지표(각 계약의 입차수와 받은 양의 합)와 얼마나 잘 맞는지 나타낸 분야 일치도예요. 같은 기간에서 쟀어요(매칭 코호트, 계약 $n=6{,}483$개). 점수의 순위와 대리 지표의 순위 사이의 켄달의 $\tau$와 스피어먼의 $\rho$를 보여 줘요. 신뢰구간(CI)은 계약을 다시 뽑아 구한 95\% 백분위 부트스트랩 구간이에요(짝지은 다시 뽑기 400번, 시드 42).}
\label{tab:alignment-transfer}
\fitwidth{%
\begin{tabular}{lccc}
\toprule
대리 지표 & 스피어먼의 $\rho$ & 켄달의 $\tau$ & $\tau$의 95\% 신뢰구간 \\
\midrule
\multicolumn{4}{l}{\textit{EndorseRank}} \\
입차수 & 0.302 & 0.218 & [0.199, 0.238] \\
받은 양의 합 (들어온 양을 모두 더함) & 0.256 & 0.175 & [0.157, 0.194] \\
\addlinespace
\multicolumn{4}{l}{\textit{AWP}} \\
입차수 & 0.883 & 0.742 & [0.732, 0.752] \\
받은 양의 합 (들어온 양을 모두 더함) & 0.887 & 0.730 & [0.720, 0.740] \\
\bottomrule
\end{tabular}
}
\end{table}


AWP는 입차수와 맞는 것과 거의 같은 만큼 받은 양의 합과도 잘 맞아요(표~\ref{tab:alignment-transfer}). 그 까닭은 \ref{sec:same-window-alignment}절에 있어요. 엔도스랭크의 송금 계수의 구간은 0과도 겹치지 않고, AWP의 구간과도 겹치지 않아요. 이 계수들은 주로, 많은 주인이 허락한 계약이 코호트의 나머지 계약보다 더 많은 보낸 쪽(송신자)에게서 송금을 받는다는 것을 보여 줘요.

\dissertationSubsubheading[sub:allowance-family]{허락 가족}

표~\ref{tab:alignment-allowance}은 두 점수와 받은 허락 수, 그리고 받은 허락 금액 사이의 계수를 보여 줘요.

% Generated by experiments/scripts/export_latex.py -- do not edit by hand.
% Source: experiments/data/2-processed-tables-and-evaluations/revised-usd-12m/same-window/eval_summary.json
\begin{table}[htbp]
\centering
\caption[허락 대리 지표에 대한 엔도스랭크와 AWP의 분야 일치도.]{엔도스랭크와 AWP가 허락 대리 지표(각 스펜더 계약의 받은 허락 수와 받은 허락의 양)와 얼마나 잘 맞는지 나타낸 분야 일치도예요. 같은 기간에서 쟀어요(매칭 코호트, 계약 $n=6{,}483$개). 점수의 순위와 대리 지표의 순위 사이의 켄달의 $\tau$와 스피어먼의 $\rho$를 보여 줘요. 신뢰구간(CI)은 계약을 다시 뽑아 구한 95\% 백분위 부트스트랩 구간이에요(짝지은 다시 뽑기 400번, 시드 42).}
\label{tab:alignment-allowance}
\fitwidth{%
\begin{tabular}{lccc}
\toprule
대리 지표 & 스피어먼의 $\rho$ & 켄달의 $\tau$ & $\tau$의 95\% 신뢰구간 \\
\midrule
\multicolumn{4}{l}{\textit{EndorseRank}} \\
받은 허락 수 & 0.804 & 0.650 & [0.639, 0.661] \\
받은 허락의 양 (합) & 0.628 & 0.469 & [0.454, 0.484] \\
\addlinespace
\multicolumn{4}{l}{\textit{AWP}} \\
받은 허락 수 & 0.320 & 0.242 & [0.223, 0.262] \\
받은 허락의 양 (합) & 0.294 & 0.211 & [0.193, 0.229] \\
\bottomrule
\end{tabular}
}
\end{table}


엔도스랭크는 받은 허락 금액보다 받은 허락 수를 더 가깝게 따라가요. 그 까닭은 \ref{sec:same-window-alignment}절에 있어요. AWP의 계수는 0보다 위에 있어요. 코호트의 계약 가운데, 많은 주인이 허락한 계약은 송금 그래프에서도 중심에 있는 경향이 있어요.

\dissertationSubsubheading[sub:sybil-family]{시빌 보정 안정성 가족}

표~\ref{tab:alignment-sybil-stability}는 두 점수와 들어오는 상대 비율, 송금 활동 기간, 활동한 달 수 사이의 계수를 보여 줘요.

% Generated by experiments/scripts/export_latex.py -- do not edit by hand.
% Source: experiments/data/2-processed-tables-and-evaluations/revised-usd-12m/same-window/eval_summary.json
\begin{table}[htbp]
\centering
\caption[시빌 보정 안정성 대리 지표에 대한 엔도스랭크와 AWP의 분야 일치도.]{엔도스랭크와 AWP가 송금 활동에서 구한 시빌 보정 안정성 대리 지표와 얼마나 잘 맞는지 나타낸 분야 일치도예요. 같은 기간에서 쟀어요(매칭 코호트, 계약 $n=6{,}483$개). 점수의 순위와 대리 지표의 순위 사이의 켄달의 $\tau$와 스피어먼의 $\rho$를 보여 줘요. 신뢰구간(CI)은 계약을 다시 뽑아 구한 95\% 백분위 부트스트랩 구간이에요(짝지은 다시 뽑기 400번, 시드 42).}
\label{tab:alignment-sybil-stability}
\fitwidth{%
\begin{tabular}{lccc}
\toprule
대리 지표 & 스피어먼의 $\rho$ & 켄달의 $\tau$ & $\tau$의 95\% 신뢰구간 \\
\midrule
\multicolumn{4}{l}{\textit{EndorseRank}} \\
들어오는 상대 비율 & -0.290 & -0.197 & [-0.213, -0.180] \\
송금 활동 기간 (일) & 0.199 & 0.133 & [0.114, 0.153] \\
활동한 달 수 & 0.203 & 0.143 & [0.124, 0.164] \\
\addlinespace
\multicolumn{4}{l}{\textit{AWP}} \\
들어오는 상대 비율 & 0.163 & 0.055 & [0.033, 0.079] \\
송금 활동 기간 (일) & 0.763 & 0.588 & [0.575, 0.600] \\
활동한 달 수 & 0.768 & 0.616 & [0.603, 0.629] \\
\bottomrule
\end{tabular}
}
\end{table}


AWP는 송금 활동 기간과 활동한 달 수에서 엔도스랭크보다 앞서요. 기간 전체에 퍼진 활동에 점수를 더 주는 송금 점수라면 그럴 것이라고 예상할 수 있어요. 지표 가족 평균에서의 짝지은 차이는 0을 벗어나요(표~\ref{tab:tau-diff}, 대비~iv). 들어오는 상대 비율에서는 두 계수의 부호가 서로 반대예요. 엔도스랭크는 음수이고, AWP는 양수예요. 이 비율은 같은 보낸 쪽이 다시 오면 내려가요. 그래서 많은 주인이 허락한 계약은, 이 잣대로 보면 보낸 쪽이 다시 찾아오는 계약이에요. AWP의 계수가 양수인 것은 점수를 매기는 범위 때문이에요. 송금 그래프 밖에 있는 계약 $1{,}192$개는 비율도 0이고 AWP 점수도 0이에요. 그래서 이 계약들이 두 줄 세우기의 맨 아래를 함께 채워요. 이 가족은 각 점수가 시간에 걸친 송금 정보를 얼마나 담고 있는지를 보여 줘요.

\dissertationSubheading[sec:er-awp-holdout]{기간 밖}

표~\ref{tab:endorserank-awp-holdout}은 두 홀드아웃(나중 기간을 떼어 두고 맞혀 보는 시험) 기간에서 두 점수 사이의 짝지은 차이를 모아요. 모든 행에서 $a$는 엔도스랭크이고 $b$는 AWP예요. 스펜더(spender, 허락을 받아 대신 쓰는 쪽) 행은 라벨을 본 뒤에 계산했어요.

% Generated by experiments/scripts/export_latex.py -- do not edit by hand.
% Source: experiments/data/2-processed-tables-and-evaluations/revised-usd-12m/holdout-w0/eval_summary.json
% Source: experiments/data/2-processed-tables-and-evaluations/revised-usd-12m/supplementary/supplementary.json
% Source: experiments/data/2-processed-tables-and-evaluations/revised-usd-12m/window-w1/eval_summary.json
\begin{table}[htbp]
\centering
\caption[기간 밖에서 엔도스랭크와 AWP 사이의 짝지은 부트스트랩 차이.]{기간 밖에서 엔도스랭크와 AWP 사이의 차이를 짝지은 부트스트랩으로 구했어요(W0 스펜더 $n=7{,}914$, W0 트레이더 $n=87$, W1 스펜더 $n=7{,}960$, W1 트레이더 $n=70$). 이 가운데 어느 것도 규칙 A나 규칙 B에 들어가지 않아요. 스펜더 행은 라벨을 본 뒤에 계산했어요. 트레이더 행은 어느 화살표 종류로도 만들지 않은 결과에서 한 비교 가운데, 엔도스랭크에서 AWP를 뺀 쌍이에요(표~\ref{tab:neutral-label-contrasts}와~\ref{tab:neutral-label-grid}). W1 스펜더 행은 W1 평가가 끝난 뒤에, 저장해 둔 W1 점수로 계산했어요. 다시 뽑은 표본은 표~\ref{tab:fresh-holdout}와 같은 것을 썼어요. 각 묶음 안에서 $a$와 $b$는 같은 계약 다시 뽑은 표본을 함께 써요(스펜더는 400번, 트레이더는 2{,}000번, 시드 42).}
\label{tab:endorserank-awp-holdout}
\fitwidth{%
\begin{tabular}{llccccc}
\toprule
대비 ($a$ 빼기 $b$) & 라벨 & $\tau(a)$ & $\tau(b)$ & $\Delta\tau$ & $\Delta\tau$의 95\% 신뢰구간 & $P(\Delta\tau>0)$ \\
\midrule
\multicolumn{7}{l}{\emph{W0, 스펜더 ($n=7{,}914$)}} \\
EndorseRank 빼기 AWP & 새 허락 쌍 & 0.303 & 0.196 & +0.107 & [0.086, 0.130] & 1.000 \\
EndorseRank 빼기 AWP & 새 송금자 & 0.272 & 0.409 & -0.138 & [-0.151, -0.123] & 0.000 \\
\midrule
\multicolumn{7}{l}{\emph{W0, 세 번 이상 닫은 트레이더 ($n=87$)}} \\
EndorseRank 빼기 AWP & 청산 없이 닫은 비율 & 0.236 & 0.254 & -0.019 & [-0.173, 0.144] & 0.401 \\
\midrule
\multicolumn{7}{l}{\emph{W1, 스펜더 ($n=7{,}960$)}} \\
EndorseRank 빼기 AWP & 새 허락 쌍 & 0.299 & 0.157 & +0.142 & [0.119, 0.164] & 1.000 \\
EndorseRank 빼기 AWP & 새 송금자 & 0.259 & 0.398 & -0.139 & [-0.154, -0.125] & 0.000 \\
\midrule
\multicolumn{7}{l}{\emph{W1, 세 번 이상 닫은 트레이더 ($n=70$)}} \\
EndorseRank 빼기 AWP & 청산 없이 닫은 비율 & 0.297 & 0.263 & +0.034 & [-0.067, 0.134] & 0.750 \\
\bottomrule
\end{tabular}
}
\end{table}


W0의 스펜더에서 엔도스랭크는 새 허락 쌍을 AWP보다 더 잘 줄 세워요. 짝지은 차이는 $+0.107$(구간 $[0.086,0.130]$)이에요. 새 송금자에서는 순서가 뒤집혀요. 차이는 $-0.138$($[-0.151,-0.123]$)이에요. 모든 라벨은 점수에 쓴 모든 이벤트보다 뒤에 와요. 첫 번째 차이는 점수를 매기는 범위에 달려 있어요. AWP는 동결일 전에 송금 화살표가 없던 스펜더 $1{,}342$개에 0점을 줘요. 이것은 대략 여섯 개에 하나꼴이에요. 이 스펜더들은 허락 화살표로만 순서가 정해져요. 그리고 이 스펜더들은 나머지 스펜더보다 새 허락 쌍을 더 자주 얻었어요($26.7\%$ 대 $20.9\%$). 다른 주소에게서 송금을 받은 적이 있는 스펜더 $6{,}489$개에서는 AWP가 두 라벨 모두에서 앞서요(표~\ref{tab:holdout-receiving}). 여기에는 짝지은 검사를 하지 않았어요.

7월부터 9월까지에서도 AWP는 송금 화살표가 없는 스펜더에 다시 0점을 줘요. 그리고 같은 모양이 되풀이돼요. 엔도스랭크는 새 허락 쌍에서 앞서고($+0.142$, $[0.119,0.164]$), AWP는 새 송금자에서 앞서요($-0.139$, $[-0.154,-0.125]$).

W0 트레이더에서는(표~\ref{tab:holdout-traders}) 이익을 낸 닫기의 수와 실현 이익(실제로 얻은 이익)에서 AWP의 구간만 0보다 위에 있어요. 이 두 라벨은 닫기 횟수가 많을수록 커져요. 두 비율은 그렇지 않은데, 여기서는 두 점수가 모두 음수예요. 그리고 손실이 아닌 비율에서만 구간이 0보다 아래에 있어요. 어느 점수도 이 계약 트레이더(GMX~V2에서 거래하는 계약)들을 성공한 닫기의 비율대로 줄 세우지 못해요.

\dissertationSubheading[sec:er-awp-neutral-grid]{두 기간의 모든 트레이더 라벨}

표~\ref{tab:neutral-label-grid}은 두 기간에서 트레이더 라벨 여섯 개 모두에 대해 허락 쪽과 송금 쪽을 맞세워요. 이 표의 쌍은 이 연구가 7월부터 9월까지의 로그를 뽑기 전에 W1과 함께 등록했어요(\ref{sub:neutral-label-protocol}절). 여기서는 등록 분석의 결과를 안 뒤에 그 쌍을 고친 점수로 계산해요. 그래서 이 표는 재분석이에요(부록~\ref{sec:map-registered}). 청산 없이 닫은 비율에서의 P와 Q는 규칙 R1과 R2에 들어가요. 그리고 두 비율에서의 P는 규칙 R3에 들어가요(주장~C6, \ref{sub:neutral-label-results}절). W1에서는 두 청산 라벨 모두에서 개수(P)가 허락 쪽에 유리하고, 구간이 0보다 위에 있어요. 반면 세 페이지랭크 쌍은 양수이지만, 구간이 0을 품어요. W0에서는 두 청산 라벨 모두에서 P가 양수이지만, 구간이 0을 품어요. 그리고 세 페이지랭크 쌍은 0보다 아래에 있고, 그 구간도 0을 품어요. W1의 이익 라벨 네 개에서는 모든 쌍이 송금 쪽에 유리해요. Q, C-PR 층, 엔도스랭크 쌍의 구간은 네 라벨 모두에서 0보다 아래에 있어요. P의 구간은 두 비율에서 0보다 아래에 있어요. W0에서는 이익 라벨의 모든 구간이 0을 품어요. 다만 실현 이익에서의 C-PR 층의 구간은 예외예요.

% Generated by experiments/scripts/export_latex.py -- do not edit by hand.
% Source: experiments/data/2-processed-tables-and-evaluations/revised-usd-12m/window-w1/eval_summary.json
\begin{table}[htbp]
\centering
\caption[두 기간의 모든 트레이더 라벨에서 허락 쪽 점수 빼기 송금 쪽 점수.]{허락 쪽 점수에서 송금 쪽 점수를 뺀 $\Delta\tau$와 95\% 짝지은 부트스트랩 구간이에요. 기간 W1(GMX~V2 계약 트레이더 $n=70$)과 W0($n=87$)의 트레이더 라벨 여섯 개 모두에서 구했어요. P는 받은 허락 수 빼기 송금 입차수예요. Q는 활동 재시작을 쓰는 엔도스랭크 빼기 AWP예요(화살표 무게와 다시 출발 규칙이 같음). C-PR은 $\lambda=1$ 빼기 $\lambda=0$이에요(같은 달러 방식으로 만듦). ER은 엔도스랭크 빼기 AWP예요. 이익을 낸 닫기 수와 실현 이익은 닫은 횟수가 많을수록 커져요. 다른 네 라벨은 그렇지 않아요. 신뢰구간(CI)은 계약을 다시 뽑아 구한 95\% 백분위 부트스트랩 구간이에요(짝지은 다시 뽑기 2{,}000번, 시드 42).}
\label{tab:neutral-label-grid}
\fitwidth{%
\begin{tabular}{lcccc}
\toprule
라벨 & P (개수) & Q (페이지랭크) & C-PR 층 & ER 빼기 AWP \\
\midrule
\multicolumn{5}{l}{\emph{기간 W1, 2026년 7월--9월 (트레이더 $n=70$)}} \\
청산 없이 닫은 비율 & \begin{tabular}[c]{@{}c@{}}0.141\\[-1pt]{\footnotesize [0.037, 0.240]}\end{tabular} & \begin{tabular}[c]{@{}c@{}}0.023\\[-1pt]{\footnotesize [-0.062, 0.113]}\end{tabular} & \begin{tabular}[c]{@{}c@{}}0.019\\[-1pt]{\footnotesize [-0.093, 0.128]}\end{tabular} & \begin{tabular}[c]{@{}c@{}}0.034\\[-1pt]{\footnotesize [-0.067, 0.134]}\end{tabular} \\
청산 없음 & \begin{tabular}[c]{@{}c@{}}0.141\\[-1pt]{\footnotesize [0.044, 0.232]}\end{tabular} & \begin{tabular}[c]{@{}c@{}}0.027\\[-1pt]{\footnotesize [-0.066, 0.129]}\end{tabular} & \begin{tabular}[c]{@{}c@{}}0.021\\[-1pt]{\footnotesize [-0.087, 0.132]}\end{tabular} & \begin{tabular}[c]{@{}c@{}}0.038\\[-1pt]{\footnotesize [-0.069, 0.149]}\end{tabular} \\
이익을 낸 닫기 수 & \begin{tabular}[c]{@{}c@{}}-0.115\\[-1pt]{\footnotesize [-0.257, 0.025]}\end{tabular} & \begin{tabular}[c]{@{}c@{}}-0.215\\[-1pt]{\footnotesize [-0.399, -0.042]}\end{tabular} & \begin{tabular}[c]{@{}c@{}}-0.213\\[-1pt]{\footnotesize [-0.402, -0.028]}\end{tabular} & \begin{tabular}[c]{@{}c@{}}-0.200\\[-1pt]{\footnotesize [-0.390, -0.023]}\end{tabular} \\
실현 이익 & \begin{tabular}[c]{@{}c@{}}-0.109\\[-1pt]{\footnotesize [-0.247, 0.032]}\end{tabular} & \begin{tabular}[c]{@{}c@{}}-0.299\\[-1pt]{\footnotesize [-0.457, -0.146]}\end{tabular} & \begin{tabular}[c]{@{}c@{}}-0.294\\[-1pt]{\footnotesize [-0.476, -0.116]}\end{tabular} & \begin{tabular}[c]{@{}c@{}}-0.292\\[-1pt]{\footnotesize [-0.460, -0.134]}\end{tabular} \\
이익을 낸 비율 & \begin{tabular}[c]{@{}c@{}}-0.162\\[-1pt]{\footnotesize [-0.295, -0.025]}\end{tabular} & \begin{tabular}[c]{@{}c@{}}-0.262\\[-1pt]{\footnotesize [-0.436, -0.100]}\end{tabular} & \begin{tabular}[c]{@{}c@{}}-0.266\\[-1pt]{\footnotesize [-0.461, -0.070]}\end{tabular} & \begin{tabular}[c]{@{}c@{}}-0.246\\[-1pt]{\footnotesize [-0.427, -0.073]}\end{tabular} \\
손실이 아닌 비율 & \begin{tabular}[c]{@{}c@{}}-0.222\\[-1pt]{\footnotesize [-0.328, -0.106]}\end{tabular} & \begin{tabular}[c]{@{}c@{}}-0.324\\[-1pt]{\footnotesize [-0.471, -0.188]}\end{tabular} & \begin{tabular}[c]{@{}c@{}}-0.316\\[-1pt]{\footnotesize [-0.481, -0.152]}\end{tabular} & \begin{tabular}[c]{@{}c@{}}-0.315\\[-1pt]{\footnotesize [-0.468, -0.176]}\end{tabular} \\
\midrule
\multicolumn{5}{l}{\emph{기간 W0, 2026년 4월--6월 (트레이더 $n=87$)}} \\
청산 없이 닫은 비율 & \begin{tabular}[c]{@{}c@{}}0.095\\[-1pt]{\footnotesize [-0.027, 0.234]}\end{tabular} & \begin{tabular}[c]{@{}c@{}}-0.012\\[-1pt]{\footnotesize [-0.153, 0.149]}\end{tabular} & \begin{tabular}[c]{@{}c@{}}-0.057\\[-1pt]{\footnotesize [-0.209, 0.092]}\end{tabular} & \begin{tabular}[c]{@{}c@{}}-0.019\\[-1pt]{\footnotesize [-0.173, 0.144]}\end{tabular} \\
청산 없음 & \begin{tabular}[c]{@{}c@{}}0.109\\[-1pt]{\footnotesize [-0.022, 0.257]}\end{tabular} & \begin{tabular}[c]{@{}c@{}}-0.005\\[-1pt]{\footnotesize [-0.162, 0.165]}\end{tabular} & \begin{tabular}[c]{@{}c@{}}-0.056\\[-1pt]{\footnotesize [-0.211, 0.098]}\end{tabular} & \begin{tabular}[c]{@{}c@{}}-0.007\\[-1pt]{\footnotesize [-0.170, 0.165]}\end{tabular} \\
이익을 낸 닫기 수 & \begin{tabular}[c]{@{}c@{}}-0.002\\[-1pt]{\footnotesize [-0.141, 0.129]}\end{tabular} & \begin{tabular}[c]{@{}c@{}}-0.027\\[-1pt]{\footnotesize [-0.170, 0.113]}\end{tabular} & \begin{tabular}[c]{@{}c@{}}-0.089\\[-1pt]{\footnotesize [-0.244, 0.070]}\end{tabular} & \begin{tabular}[c]{@{}c@{}}-0.046\\[-1pt]{\footnotesize [-0.205, 0.113]}\end{tabular} \\
실현 이익 & \begin{tabular}[c]{@{}c@{}}-0.026\\[-1pt]{\footnotesize [-0.176, 0.112]}\end{tabular} & \begin{tabular}[c]{@{}c@{}}-0.110\\[-1pt]{\footnotesize [-0.270, 0.040]}\end{tabular} & \begin{tabular}[c]{@{}c@{}}-0.257\\[-1pt]{\footnotesize [-0.447, -0.071]}\end{tabular} & \begin{tabular}[c]{@{}c@{}}-0.139\\[-1pt]{\footnotesize [-0.311, 0.022]}\end{tabular} \\
이익을 낸 비율 & \begin{tabular}[c]{@{}c@{}}-0.043\\[-1pt]{\footnotesize [-0.182, 0.081]}\end{tabular} & \begin{tabular}[c]{@{}c@{}}0.044\\[-1pt]{\footnotesize [-0.084, 0.183]}\end{tabular} & \begin{tabular}[c]{@{}c@{}}0.012\\[-1pt]{\footnotesize [-0.125, 0.150]}\end{tabular} & \begin{tabular}[c]{@{}c@{}}0.074\\[-1pt]{\footnotesize [-0.066, 0.223]}\end{tabular} \\
손실이 아닌 비율 & \begin{tabular}[c]{@{}c@{}}-0.074\\[-1pt]{\footnotesize [-0.208, 0.052]}\end{tabular} & \begin{tabular}[c]{@{}c@{}}0.078\\[-1pt]{\footnotesize [-0.055, 0.217]}\end{tabular} & \begin{tabular}[c]{@{}c@{}}0.047\\[-1pt]{\footnotesize [-0.090, 0.191]}\end{tabular} & \begin{tabular}[c]{@{}c@{}}0.108\\[-1pt]{\footnotesize [-0.034, 0.256]}\end{tabular} \\
\bottomrule
\end{tabular}
}
\end{table}

